% ============================================================
% Section 104: AB原子线的一维X射线衍射
% ============================================================
\section{固体理论：AB原子线的一维X射线衍射}
\label{sec:ab-chain-diffraction}

\begin{modelbox}[题目：AB原子线衍射条件、强度与系统消光]
考虑一个 $ABAB\cdots AB$ 原子线，$A-B$ 键长为 $a/2$。$A$、$B$ 原子的散射因子分别为 $f_A$ 和 $f_B$，入射 X 射线垂直于原子线。

(1) 给出 $\theta$ 方向（衍射光束与原子线夹角）衍射加强条件；
(2) 计算衍射强度；
(3) 讨论 $f_A=f_B$ 的情况。
\end{modelbox}

\begin{tipbox}[考点快速分析]
\begin{itemize}
    \item \textbf{一维衍射条件}：$a\cos\theta=n\lambda$（原胞长度 $a$，含一个 $A$ 和一个 $B$）
    \item \textbf{结构因子}：$F=f_A+f_Be^{i\pi n}=f_A+(-1)^nf_B$
    \item \textbf{强度}：$I\propto|F|^2=|f_A+(-1)^nf_B|^2$
    \item \textbf{系统消光}：$f_A=f_B$ 时奇数级 $F=0$，完全消光
\end{itemize}
\end{tipbox}

\begin{derivationbox}[标准解题过程]
\textbf{(1) 衍射加强条件}

原子线可视为一维晶体，原胞长度为 $a$（每个原胞含一个 $A$ 和一个 $B$，相距 $a/2$）。X 射线垂直于原子线入射，衍射光束与原子线夹角为 $\theta$。

相邻原胞散射波的程差为 $a\cos\theta$。衍射加强条件为
\begin{equation}
\boxed{a\cos\theta=n\lambda,\qquad n=0,\pm1,\pm2,\ldots}
\end{equation}

\textbf{(2) 衍射强度的计算}

衍射强度正比于晶体总散射振幅的模平方。总振幅由原胞结构因子 $F$ 决定：
\begin{equation}
I\propto N^2|F|^2,
\end{equation}
其中 $N$ 为原胞数。

取 $A$ 原子为原点（$x_A=0$），$B$ 原子位于 $x_B=a/2$。倒格矢 $G=2\pi n/a$。几何结构因子
\begin{equation}
F=f_Ae^{iG\cdot0}+f_Be^{iG\cdot a/2}
=f_A+f_Be^{i\pi n}
=f_A+(-1)^nf_B.
\end{equation}

因此衍射强度
\begin{equation}
\boxed{I\propto|f_A+(-1)^nf_B|^2=}
\begin{cases}
|f_A+f_B|^2, & n\text{ 为偶数},\\
|f_A-f_B|^2, & n\text{ 为奇数}.
\end{cases}
\end{equation}

\textbf{(3) $f_A=f_B$ 的情况}

令 $f_A=f_B=f$（如 KCl 中 K$^+$ 与 Cl$^-$ 电子数相近），则
\begin{equation}
F=\begin{cases}2f, & n\text{ 为偶数},\\0, & n\text{ 为奇数}.\end{cases}
\end{equation}

衍射强度
\begin{equation}
\boxed{I\propto\begin{cases}4f^2, & n\text{ 为偶数},\\0, & n\text{ 为奇数}.\end{cases}}
\end{equation}

\textbf{物理意义}：
\begin{itemize}
    \item $n$ 为偶数时，$A$、$B$ 原子散射波相位差为 $2\pi$ 的整数倍，同相叠加，强度最大。
    \item $n$ 为奇数时，两原子散射波相位差为 $\pi$，振幅完全抵消，发生\textbf{系统消光}。
\end{itemize}
\end{derivationbox}

\begin{formulabox}[答案汇总]
\begin{center}
\begin{tabular}{ll}
\toprule
\textbf{项目} & \textbf{结果} \\
\midrule
衍射条件 & $a\cos\theta=n\lambda$ \\
衍射强度 & $I\propto|f_A+(-1)^nf_B|^2$ \\
$f_A=f_B$ 时 & 偶数级 $I\propto4f^2$（强）；奇数级 $I=0$（系统消光） \\
\bottomrule
\end{tabular}
\end{center}
\end{formulabox}

\begin{insightbox}[物理图像]
\begin{itemize}
    \item 一维晶格的衍射条件由原胞间距 $a$ 决定，与三维布拉格定律在形式上有所不同（三维为 $2d\sin\theta=n\lambda$）。
    \item 衍射强度由几何结构因子 $F$ 决定，$F$ 反映了原胞内原子的排列和散射能力。当 $F=0$ 时，即使满足衍射条件，也无衍射线出现——这就是\textbf{系统消光}。
    \item 系统消光是晶体结构分析的重要判据，可用于推断原胞内原子的相对位置和种类。AB 链的消光规律与 NaCl、金刚石等三维复式格子的消光规律同源。
\end{itemize}
\end{insightbox}
